\begin{center}
    {\LARGE \textbf{Parte III}}\\[6pt]
\end{center}
\vspace{4pt}\hrule\vspace{10pt}

% ==================================================================
\subsection*{Entendiendo el \textit{environment}}
% ==================================================================

Se considera una economía mundial formada por $N \geq 2$ países indexados por
$i \in \{1, \ldots, N\}$.
Cada país $i$ tiene preferencias \textbf{logarítmicas} con factores de descuento
\textbf{heterogéneos}:
\[
    U_i = \ln(c_{i,1}) + \beta_i \ln(c_{i,2}), \qquad \beta_i \in (0,1).
\]

La \textbf{restricción presupuestaria} del hogar representativo en cada período es:
\begin{equation}\label{eq:rp}
    c_{i,t} + (1+r^*)\,d_{i,t-1} = y_{i,t} + d_{i,t},
\end{equation}
donde $d_{i,t}>0$ indica deuda neta asumida en $t$ y exigible en $t+1$,
y $r^*>0$ es la tasa de interés mundial (endógena en esta parte).

La condición de \textbf{no-Ponzi} es:
\[
    \lim_{j\to\infty} \mathbb{E}_t \frac{d_{i,t+j}}{(1+r^*)^j} \leq 0.
\]

Se asume que $\beta_i(1+r^*)=1$ y que cada hogar \textbf{inicia sin deuda}:
$d_{i,0}=0$.

Se definen la \textbf{balanza comercial} y la \textbf{cuenta corriente} como:
\[
    tb_{i,t} = y_{i,t} - c_{i,t}, \qquad
    ca_{i,t} = tb_{i,t} - r^*\,d_{i,t-1} = -(d_{i,t} - d_{i,t-1}).
\]

% ==================================================================
\subsection*{Ítem 17}
% ==================================================================

\subsubsection*{Planteamiento}

El hogar del país $i$ resuelve:
\[
    \max_{c_{i,1},\, c_{i,2}} \;\; \ln(c_{i,1}) + \beta_i \ln(c_{i,2}).
\]

\paragraph{Variables de control y de estado.}
\begin{itemize}
    \item \textbf{Variables de control:} $c_{i,1}$, $c_{i,2}$.
    \item \textbf{Variables de estado:} $\lambda$ (multiplicador de Lagrange), $r^*$.
\end{itemize}

\subsubsection*{Derivación de la Restricción Presupuestaria Intertemporal (RPI)}

Escribimos las restricciones período a período usando \eqref{eq:rp}.
Con $d_{i,0}=0$ y la condición de transversalidad ($d_{i,2}=0$, los hogares
mueren sin deuda):

\medskip
\textit{Período 1:}
\[
    c_{i,1} + \underbrace{(1+r^*)\cdot 0}_{=0} = y_{i,1} + d_{i,1}
    \quad\Longrightarrow\quad
    c_{i,1} = y_{i,1} + d_{i,1}.
\]

\textit{Período 2:}
\[
    c_{i,2} + (1+r^*)\,d_{i,1} = y_{i,2} + \underbrace{d_{i,2}}_{=0}
    \quad\Longrightarrow\quad
    c_{i,2} + (1+r^*)\,d_{i,1} = y_{i,2}.
\]

Despejamos $d_{i,1}$ de la ecuación del período 1:
\[
    d_{i,1} = c_{i,1} - y_{i,1}.
\]
Sustituimos en la ecuación del período 2:
\[
    c_{i,2} + (1+r^*)(c_{i,1} - y_{i,1}) = y_{i,2}.
\]
Expandimos y reordenamos llevando los términos de consumo a la izquierda y los
de ingreso a la derecha:
\[
    c_{i,2} + (1+r^*)c_{i,1} - (1+r^*)y_{i,1} = y_{i,2}.
\]
\[
    (1+r^*)c_{i,1} + c_{i,2} = y_{i,2} + (1+r^*)y_{i,1}.
\]
Dividimos todo entre $(1+r^*)$:

\begin{equation}\label{eq:rpi}
    \boxed{c_{i,1} + \frac{c_{i,2}}{1+r^*} = y_{i,1} + \frac{y_{i,2}}{1+r^*}.}
\end{equation}

Esta es la \textbf{RPI}: el valor presente del consumo debe igualar el valor
presente del ingreso.

\subsubsection*{Lagrangiano y Condiciones de Primer Orden (CPO)}

Formamos el Lagrangiano:
\[
\mathcal{L} = \ln(c_{i,1}) + \beta_i \ln(c_{i,2})
              + \lambda \!\left[y_{i,1} + \frac{y_{i,2}}{1+r^*}
              - c_{i,1} - \frac{c_{i,2}}{1+r^*}\right].
\]

\paragraph{CPO respecto a $c_{i,1}$:}
\[
    \frac{\partial \mathcal{L}}{\partial c_{i,1}} = \frac{1}{c_{i,1}} - \lambda = 0
    \quad\xrightarrow{(b)}\quad
    \lambda = \frac{1}{c_{i,1}}.
    \tag{b}
\]

\paragraph{CPO respecto a $c_{i,2}$:}
\[
    \frac{\partial \mathcal{L}}{\partial c_{i,2}}
    = \frac{\beta_i}{c_{i,2}} - \frac{\lambda}{1+r^*} = 0
    \quad\xrightarrow{(b)}\quad
    \lambda = \frac{(1+r^*)\beta_i}{c_{i,2}}.
    \tag{b}
\]

\paragraph{CPO respecto a $\lambda$ (factibilidad):}
\[
    \frac{\partial \mathcal{L}}{\partial \lambda}
    = y_{i,1} + \frac{y_{i,2}}{1+r^*} - c_{i,1} - \frac{c_{i,2}}{1+r^*} = 0.
    \tag{iv}
\]

\subsubsection*{Ecuación de Euler}

Igualamos las dos expresiones para $\lambda$:
\[
    \frac{1}{c_{i,1}} = \frac{(1+r^*)\beta_i}{c_{i,2}}.
\]
Multiplicamos cruzado:
\[
    c_{i,2} = (1+r^*)\,\beta_i\, c_{i,1}.
    \tag{5}\label{eq:euler}
\]
Esta es la \textbf{ecuación de Euler}: el consumo crece de período 1 a período 2
a una tasa que depende de la paciencia $\beta_i$ y de la tasa de interés.

\subsubsection*{Condición de transversalidad y condición de limpieza}

\begin{condicion}[Transversalidad]
Los hogares mueren sin deuda al final del período 2:
\[
    d_{i,2} = 0.
\]
\textit{Intuición: ``No tiene sentido guardar todo y nunca disfrutarlo.''}
\end{condicion}

\begin{condicion}[Limpieza de mercado --- suma de deudas]
En el mercado internacional de deuda, la posición del mundo en el modelo
SOE debe ser cero. La suma de las deudas de todos los países debe ser cero:
\[
    \sum_{i=1}^{N} d_{i,1} = 0.
\]
\textit{Hay alguien que presta y otro que recibe, la suma debe ser cero.}
\end{condicion}

% ==================================================================
\subsection*{Ítem 18}
% ==================================================================

\subsubsection*{Consumo óptimo en período 1}

Partimos de la RPI \eqref{eq:rpi}:
\[
    c_{i,1} + \frac{c_{i,2}}{1+r^*} = y_{i,1} + \frac{y_{i,2}}{1+r^*}.
\]

Sustituimos la ecuación de Euler \eqref{eq:euler},
$c_{i,2} = (1+r^*)\beta_i c_{i,1}$, en la RPI:
\[
    c_{i,1} + \frac{(1+r^*)\beta_i\, c_{i,1}}{1+r^*}
    = y_{i,1} + \frac{y_{i,2}}{1+r^*}.
\]
Simplificamos $(1+r^*)$ en el segundo término del lado izquierdo:
\[
    c_{i,1} + \beta_i\, c_{i,1} = y_{i,1} + \frac{y_{i,2}}{1+r^*}.
\]
Factorizamos $c_{i,1}$:
\[
    c_{i,1}(1 + \beta_i) = y_{i,1} + \frac{y_{i,2}}{1+r^*}.
\]
Despejamos:
\begin{equation}\label{eq:c1}
    \boxed{c_{i,1}^* = \frac{1}{1+\beta_i}\left[y_{i,1} + \frac{y_{i,2}}{1+r^*}\right].}
    \tag{6}
\end{equation}

\subsubsection*{Consumo óptimo en período 2}

Sustituimos $c_{i,1}^*$ en la ecuación de Euler \eqref{eq:euler}:
\[
    c_{i,2}^* = (1+r^*)\,\beta_i\, c_{i,1}^*
    = (1+r^*)\,\beta_i \cdot \frac{1}{1+\beta_i}
      \left[y_{i,1} + \frac{y_{i,2}}{1+r^*}\right].
\]

\begin{equation}\label{eq:c2}
    \boxed{c_{i,2}^* = \frac{\beta_i(1+r^*)}{1+\beta_i}
    \left[y_{i,1} + \frac{y_{i,2}}{1+r^*}\right].}
    \tag{7}
\end{equation}

\subsubsection*{Deuda óptima}

\textit{Nota: solo hay $d_{i,1}$ en el primer período, pues por transversalidad
$d_{i,2}=0$.}

De la restricción presupuestaria del período 1:
\[
    c_{i,1} = y_{i,1} + d_{i,1}
    \quad\Longrightarrow\quad
    d_{i,1} = c_{i,1} - y_{i,1}.
    \tag{8}
\]

Sustituimos $c_{i,1}^*$ de \eqref{eq:c1} en (8):
\[
    d_{i,1}^* = \frac{1}{1+\beta_i}\left[y_{i,1} + \frac{y_{i,2}}{1+r^*}\right] - y_{i,1}.
\]
Ponemos $y_{i,1}$ con el mismo denominador $(1+\beta_i)$:
\[
    d_{i,1}^* = \frac{y_{i,1} + \dfrac{y_{i,2}}{1+r^*} - (1+\beta_i)\,y_{i,1}}{1+\beta_i}.
\]
Expandimos el numerador:
\[
    d_{i,1}^* = \frac{y_{i,1} + \dfrac{y_{i,2}}{1+r^*} - y_{i,1} - \beta_i\,y_{i,1}}{1+\beta_i}
    = \frac{\dfrac{y_{i,2}}{1+r^*} - \beta_i\,y_{i,1}}{1+\beta_i}.
\]

\begin{resultado}
\[
    d_{i,1}^* = \frac{\dfrac{y_{i,2}}{1+r^*} - \beta_i\,y_{i,1}}{1+\beta_i}.
\]
\end{resultado}

\subsubsection*{¿Deudor neto o acreedor neto?}

\begin{condicion}[Deudor neto --- $d_{i,1}^* > 0$]
\[
    d_{i,1}^* > 0
    \;\Longleftrightarrow\;
    \frac{y_{i,2}}{1+r^*} > \beta_i\,y_{i,1}
    \;\Longleftrightarrow\;
    y_{i,2} > (1+r^*)\,\beta_i\,y_{i,1}.
\]
\textit{Intuición: el ingreso futuro es relativamente alto, el hogar prefiere
endeudarse hoy para suavizar consumo. La impaciencia al consumo es mayor
que el premio al ahorro $\Rightarrow$ el hogar es relativamente impaciente
a la tasa internacional.}
\end{condicion}

\begin{condicion}[Acreedor neto --- $d_{i,1}^* < 0$]
\[
    d_{i,1}^* < 0
    \;\Longleftrightarrow\;
    y_{i,2} < (1+r^*)\,\beta_i\,y_{i,1}.
\]
\textit{Intuición: $(1+r^*)\beta_i > 1$ premia el ahorro más de lo que los
hogares valoran el consumo presente $\Rightarrow$ el mercado incentiva a ser
acreedor.}
\end{condicion}

\subsubsection*{Balanza Comercial (Trade Balance)}

De la restricción presupuestaria del período 1,
$c_{i,1} = y_{i,1} + d_{i,1}$, tenemos que las exportaciones netas son:
\[
    y_{i,1} - c_{i,1} = -d_{i,1}.
\]
Por lo tanto:
\[
    TB_{i,1} = y_{i,1} - c_{i,1} = -d_{i,1}^*.
    \quad\Longrightarrow\quad \boxed{TB_{i,1} = -d_{i,1}^*.}
\]
\textit{Soy acreedor en $t=1$ (presto dinero al exterior).}

De la restricción del período 2,
$c_{i,2} + (1+r^*)d_{i,1} = y_{i,2}$:
\[
    y_{i,2} - c_{i,2} = (1+r^*)d_{i,1}.
\]
\[
    TB_{i,2} = y_{i,2} - c_{i,2} = (1+r^*)\,d_{i,1}^*.
    \quad\Longrightarrow\quad \boxed{TB_{i,2} = (1+r^*)\,d_{i,1}^*.}
\]
\textit{Soy deudor en $t=2$ (pago la deuda al exterior).}

\subsubsection*{Cuenta Corriente}

Dado que en $t=1$ el hogar no nace con deuda ($d_{i,0}=0$), la cuenta
corriente coincide con la balanza comercial:
\[
    CA_{i,t} = TB_{i,t} - r^*\,d_{i,t-1}.
\]

\textit{Período 1:} Como $d_{i,0}=0$:
\[
    CA_{i,1} = TB_{i,1} - r^*\cdot 0 = TB_{i,1} = -d_{i,1}^*.
\]

Sustituyendo la expresión de $d_{i,1}^*$:
\[
    CA_{i,1} = -d_{i,1}^* = -\left[\frac{\dfrac{y_{i,2}}{1+r^*} - \beta_i\,y_{i,1}}{1+\beta_i}\right]
    = \frac{\beta_i\,y_{i,1} - \dfrac{y_{i,2}}{1+r^*}}{1+\beta_i}.
\]

\begin{resultado}
\[
    CA_{i,1} = \frac{\beta_i\,y_{i,1} - \dfrac{y_{i,2}}{1+r^*}}{1+\beta_i}.
\]
\end{resultado}

\textit{Período 2:}
\[
    CA_{i,2} = TB_{i,2} = (1+r^*)\,d_{i,1}^*
    = (1+r^*)\left[\frac{\dfrac{y_{i,2}}{1+r^*} - \beta_i\,y_{i,1}}{1+\beta_i}\right].
\]

\begin{resultado}
\[
    CA_{i,2} = \frac{y_{i,2} - (1+r^*)\,\beta_i\,y_{i,1}}{1+\beta_i}.
\]
\end{resultado}

% ==================================================================
\subsection*{Ítem 19}
% ==================================================================

\subsubsection*{Condición de limpieza del mercado de deuda}

La condición de limpieza es $\displaystyle\sum_{i=1}^{N} d_{i,1} = 0$.
Sustituimos la expresión de $d_{i,1}^*$:

\[
    \sum_{i=1}^{N}
    \frac{\dfrac{y_{i,2}}{1+r^*} - \beta_i\,y_{i,1}}{1+\beta_i} = 0.
\]

Separamos los dos términos de la suma:
\[
    \sum_{i=1}^{N} \frac{y_{i,2}}{(1+r^*)(1+\beta_i)}
    - \sum_{i=1}^{N} \frac{\beta_i\,y_{i,1}}{1+\beta_i} = 0.
\]

Pasamos el segundo término al lado derecho:
\[
    \sum_{i=1}^{N} \frac{y_{i,2}}{(1+r^*)(1+\beta_i)}
    = \sum_{i=1}^{N} \frac{\beta_i\,y_{i,1}}{1+\beta_i}.
\]

Factorizamos $\dfrac{1}{1+r^*}$ del lado izquierdo:
\[
    \frac{1}{1+r^*}\sum_{i=1}^{N} \frac{y_{i,2}}{1+\beta_i}
    = \sum_{i=1}^{N} \frac{\beta_i\,y_{i,1}}{1+\beta_i}.
\]

Despejamos $1+r^*$:

\begin{resultado}
\begin{equation}\label{eq:rstar}
    1 + r^* = \dfrac{\displaystyle\sum_{i=1}^{N} \dfrac{y_{i,2}}{1+\beta_i}}
                    {\displaystyle\sum_{i=1}^{N} \dfrac{\beta_i\,y_{i,1}}{1+\beta_i}}.
    \tag{23}
\end{equation}
\end{resultado}

\subsubsection*{Caso: factor de descuento común ($\beta_i = \beta$ para todo $i$)}

Si todos los países tienen la misma paciencia $\beta$, el factor
$\dfrac{1}{1+\beta}$ (en numerador) y $\dfrac{\beta}{1+\beta}$ (en denominador)
son constantes y salen de las sumas:
\[
    1 + r^* = \dfrac{\dfrac{1}{1+\beta}\displaystyle\sum_{i=1}^{N} y_{i,2}}
               {\dfrac{\beta}{1+\beta}\displaystyle\sum_{i=1}^{N} y_{i,1}}.
\]
Cancelamos el factor común $(1+\beta)$ en numerador y denominador:
\[
    1 + r^* = \frac{\displaystyle\sum_{i=1}^{N} y_{i,2}}
                   {\beta \displaystyle\sum_{i=1}^{N} y_{i,1}}.
\]
Multiplicamos ambos lados por $\beta$:

\begin{resultado}
\[
    (1+r^*)\,\beta = \frac{\displaystyle\sum_{i=1}^{N} y_{i,2}}
                          {\displaystyle\sum_{i=1}^{N} y_{i,1}}
                  = \frac{\bar{Y}_2}{\bar{Y}_1},
\]
donde $\bar{Y}_t = \sum_{i=1}^N y_{i,t}$ es la \textbf{dotación mundial} en el
período $t$.
\end{resultado}

\begin{interpretacion}
Desde la perspectiva del \textbf{mercado de fondos prestables}:
\begin{itemize}
    \item El \textbf{numerador} de \eqref{eq:rstar} representa la oferta de bienes
          futuros ponderada por la paciencia de cada país.
    \item El \textbf{denominador} representa la demanda de ahorro presente
          ponderada por la paciencia.
    \item Con $\beta$ común: la tasa de interés queda determinada únicamente por
          el cociente de dotaciones mundiales $\bar{Y}_2 / \bar{Y}_1$, es decir,
          por el \textbf{crecimiento agregado de la economía global}.
    \item La ecuación $(1+r^*)\beta = \bar{Y}_2/\bar{Y}_1$ es la ecuación de Euler
          \textit{agregada}: la tasa de interés compensa exactamente la impaciencia
          común ajustada por el crecimiento mundial.
\end{itemize}
\end{interpretacion}

\subsubsection*{Descomposición de la cuenta corriente en componente agregado
e idiosincrático}

\paragraph{Definiciones.}
Definimos:
\begin{itemize}
    \item Dotación mundial promedio en $t$:
          $\bar{Y}_t \equiv \displaystyle\sum_{i=1}^N y_{i,t}$
          y $\dfrac{\bar{Y}_t}{N}$ es el promedio por país.
    \item Componente idiosincrático del ingreso del país $i$:
          \[
              y_{t,i} = \frac{\bar{Y}_t}{N} + \hat{y}_{t,i},
          \]
          donde $\hat{y}_{t,i}$ es la desviación del ingreso del país $i$
          respecto al promedio mundial.
\end{itemize}

\paragraph{Sustitución en $CA_{i,1}$.}
Partimos de:
\[
    CA_{i,1}^* = \frac{1}{1+\beta}\left(\beta\,y_{i,1} - \frac{y_{i,2}}{1+r^*}\right).
\]

Sustituimos $y_{t,i} = \dfrac{\bar{Y}_t}{N} + \hat{y}_{t,i}$:
\[
    CA_{i,1}^* = \frac{1}{1+\beta}
    \left[\beta\left(\frac{\bar{Y}_1}{N} + \hat{y}_{i,1}\right)
    - \frac{1}{1+r^*}\left(\frac{\bar{Y}_2}{N} + \hat{y}_{i,2}\right)\right].
\]

Distribuimos:
\[
    CA_{i,1}^* = \underbrace{\frac{1}{1+\beta}
    \left[\beta\cdot\frac{\bar{Y}_1}{N} - \frac{\bar{Y}_2/N}{1+r^*}\right]}_{\text{Componente agregado}}
    + \underbrace{\frac{1}{1+\beta}
    \left[\beta\,\hat{y}_{i,1} - \frac{\hat{y}_{i,2}}{1+r^*}\right]}_{\text{Componente idiosincrático}}.
\]

\paragraph{Verificación: condición de limpieza $\sum CA_{i,1}^* = 0$.}

Sumamos sobre todos los países. Con $\beta_i = \beta$:
\[
    \sum_{i=1}^N CA_{i,1}^* = \frac{1}{1+\beta}
    \left[\beta\sum_{i=1}^N y_{i,1} - \frac{\sum_{i=1}^N y_{i,2}}{1+r^*}\right].
\]
Usando la condición de equilibrio $(1+r^*)\beta = \dfrac{\bar{Y}_2}{\bar{Y}_1}$,
es decir $\dfrac{\bar{Y}_2}{1+r^*} = \beta\,\bar{Y}_1$:
\[
    \sum_{i=1}^N CA_{i,1}^* = \frac{1}{1+\beta}
    \left[\beta\,\bar{Y}_1 - \frac{\bar{Y}_2}{1+r^*}\right]
    = \frac{1}{1+\beta}\left[\beta\,\bar{Y}_1 - \beta\,\bar{Y}_1\right] = 0.
    \quad\checkmark
\]

\paragraph{Conclusión.}

\begin{resultado}
\[
    CA_{i,1}^* = \frac{1}{1+\beta}
    \left(\beta\,\hat{y}_{i,1} - \frac{\hat{y}_{i,2}}{1+r^*}\right).
\]
La cuenta corriente \textbf{solo depende del componente idiosincrático}
del ingreso de cada país ($\hat{y}_{i,t}$), no de la distribución
idiosincrática mundial agregada.
\end{resultado}

\begin{interpretacion}
\begin{itemize}
    \item La tasa de interés $r^*$ (ecuación \eqref{eq:rstar}) depende
          únicamente de los \textbf{totales mundiales} $\sum y_{i,t}$, no de
          cómo se distribuyen entre países.
    \item Los desequilibrios individuales de cuenta corriente reflejan las
          \textbf{diferencias de ingreso relativo} de cada país respecto al
          promedio mundial.
    \item El componente agregado suma cero entre países (condición de limpieza),
          por ello $\sum_{i=1}^N CA_{i,1}^* = 0$ siempre se cumple.
\end{itemize}
\end{interpretacion}
